\documentclass[11pt]{amsart}

% Encoding and fonts
\usepackage[T1]{fontenc}
\usepackage[utf8]{inputenc}
\usepackage{lmodern}

% Math packages
\usepackage{amsmath}
\usepackage{amssymb}
\usepackage{amsfonts}
\usepackage{amsthm}
\usepackage{mathtools}

% Layout and utilities
\usepackage{geometry}
\usepackage{enumitem}
\usepackage[hidelinks]{hyperref}

\geometry{margin=1in}

\numberwithin{equation}{section}
\setcounter{tocdepth}{3}
\setcounter{secnumdepth}{3}

% ---------------------------------------------------------
% Theorem-like environments
% ---------------------------------------------------------

\theoremstyle{plain}
\newtheorem{theorem}{Theorem}[section]
\newtheorem{proposition}[theorem]{Proposition}
\newtheorem{lemma}[theorem]{Lemma}
\newtheorem{corollary}[theorem]{Corollary}
\newtheorem{conjecture}[theorem]{Conjecture}

\theoremstyle{definition}
\newtheorem{definition}[theorem]{Definition}
\newtheorem{experiment}[theorem]{Experiment}
\newtheorem{heuristic}[theorem]{Heuristic}
\newtheorem{question}[theorem]{Question}

\theoremstyle{remark}
\newtheorem{remark}[theorem]{Remark}
\newtheorem{warning}[theorem]{Warning}

% ---------------------------------------------------------
% Notation
% ---------------------------------------------------------

\DeclareMathOperator{\sgn}{sgn}
\DeclareMathOperator{\li}{li}
\newcommand{\Ms}{M_{\mathrm{s}}}
\newcommand{\Msk}{M_{\mathrm{sk}}}
\newcommand{\Ap}{\mathcal{A}_{+}}
\newcommand{\Am}{\mathcal{A}_{-}}
\newcommand{\Anet}{\mathcal{A}}
\newcommand{\eps}{\varepsilon}
\newcommand{\Pp}{P^{+}}
\newcommand{\Cc}{\mathcal{C}}
\newcommand{\Sq}{\mathcal{S}}
\newcommand{\En}{\mathcal{E}}
\newcommand{\mt}{\widetilde{m}}

\title{The area function of the Mertens excursion skeleton\\
(research notebook, version 6)}

\author{}
\date{}

\begin{document}

\maketitle

\tableofcontents

\bigskip

\begin{remark}
\label{rem:exploratory_nature}
This notebook is exploratory. Its purpose is to study the Mertens function
through its excursions between consecutive zeros, and to separate what is
forced by the local geometry of those excursions from what requires
arithmetic input. It contains exact identities, conditional statements,
heuristics and archived statements from earlier versions. Every statement
carries an explicit status; statements from version 5 that are incorrect or
superseded are kept and marked as such.
\end{remark}

\section{Introduction and conventions}
\label{sec:introduction}

The main object is the Mertens function $M$, decomposed into excursions
between consecutive zeros, together with three area curves built from a
regularized skeleton of those excursions
(Section~\ref{sec:definitions}). Sections~\ref{sec:formulas}
to~\ref{sec:criteria} record exact formulas, bounds and criteria equivalent
to the Riemann hypothesis. Sections~\ref{sec:prime_horizon}
to~\ref{sec:local_large_prime} examine how prime factors above a threshold
enter $M$ and, in Section~\ref{sec:local_large_prime}, how they enter a
single excursion.

\medskip
\noindent\textbf{Convention on RH.} Every statement carries one of the labels
\emph{[unconditional]}, \emph{[under RH]} (the Riemann hypothesis is assumed;
simple zeros are stated explicitly when also needed), or \emph{[heuristic]}
(not proved; a model). An \emph{[unconditional]} equivalence with RH does
not assume RH.

\medskip
\noindent\textbf{Status lines.} Each section opens with a status line that
separates standard material, literature-based material, original
observations, heuristics and conjectures, as required by the authoring rules
of this notebook.

\medskip
\noindent\textbf{Notation changes with respect to version 5.} The symbol
$m(x)$ of version 5 is replaced by $M(x)$. The unregularized skeleton, called
$L$ in version 5, is denoted $\Msk$. The squarefree counting function, called
$Q(y)$ in version 5, is denoted $\Sq(y)$, since $Q_\nu$ denotes the quadratic
energy of an excursion. The Dirichlet energy, called $E(x)$ in version 5, is
denoted $\En(x)$. The coefficient called $A(R)$ in version 5 is denoted
$\Theta(R)$, since $\Anet$ denotes the black curve. The ratio $t/u^2$ keeps
the name $R$; the residual $M-\Msk$ of version 5, also called $R$ there, is
written out explicitly. The defect parameter of
Remark~\ref{rem:conditioning} is denoted $d_\nu$.

\section{Excursions and the regularized skeleton}
\label{sec:definitions}

\noindent\emph{Status.} Definitions and elementary facts, original to this
notebook, with complete proofs. The existence of infinitely many sign changes
of $M$ is standard; it follows for instance from the oscillation theorem of
Odlyzko and te Riele~\cite{OdlyzkoTeRiele1985}.

\medskip

Let $\mu$ be the Möbius function and
\begin{equation}
\label{eq:mertens_def}
  M(x)=\sum_{n\le x}\mu(n),\qquad x\ge 0,
\end{equation}
the Mertens function, which is piecewise constant, right-continuous, with
jumps at the integers. Let $0=z_0<z_1<z_2<\dots$ be the integers $n\ge0$ with
$M(n)=0$. Since the jumps of $M$ are $0$ or $\pm1$, every sign change of $M$
passes through a zero, and $M$ changes sign infinitely often
\emph{[unconditional]}. The $\nu$-th \emph{excursion} is
\begin{equation}
\label{eq:excursion_def}
  I_\nu=(u_\nu,v_\nu],\qquad u_\nu=z_{\nu-1},\quad v_\nu=z_\nu,
  \qquad \ell_\nu=v_\nu-u_\nu ,
\end{equation}
so that $M$ does not vanish at the integers of $(u_\nu,v_\nu)$. Inside
$I_\nu$ set
\begin{equation}
\label{eq:peak_def}
  H_\nu=\max_{u_\nu<n\le v_\nu}|M(n)|,\qquad
  \tau_\nu=\min\{n\in I_\nu:\ |M(n)|=H_\nu\},\qquad
  \sigma_\nu=\sgn M(\tau_\nu)\in\{-1,0,1\},
\end{equation}
and $a_\nu=\tau_\nu-u_\nu$, $b_\nu=v_\nu-\tau_\nu$, so that
$\ell_\nu=a_\nu+b_\nu$.

\begin{definition}[Null excursions]
\label{def:null}
An excursion is \emph{null} if $H_\nu=0$. This happens exactly when
$\ell_\nu=1$, that is, when $z_{\nu-1}$ and $z_\nu$ are consecutive integers
(equivalently $\mu(v_\nu)=0$); then $\tau_\nu=v_\nu$, $b_\nu=0$ and
$\sigma_\nu=0$. A non-null excursion has $H_\nu\ge1$ and $\ell_\nu\ge2$.
Version 5 did not distinguish null excursions; its formula for the skeleton
divides by $b_\nu$ and is undefined for them.
\end{definition}

\begin{lemma}[Height and step count]
\label{lem:height_steps}
\emph{[unconditional]} For every non-null excursion,
\begin{equation}
\label{eq:height_steps}
  H_\nu\le\min(a_\nu,b_\nu)\le \frac{\ell_\nu}{2}.
\end{equation}
\end{lemma}

\begin{proof}
The steps of $M$ are $0$ or $\pm1$, so $M$ needs at least $H_\nu$ steps to
go from $M(u_\nu)=0$ to $|M(\tau_\nu)|=H_\nu$, and at least $H_\nu$ steps to
return to $M(v_\nu)=0$.
\end{proof}

\begin{definition}[Crossing and touching zeros]
\label{def:crossing_touching}
Let $z_j<z_{j+1}<\dots<z_{j+r}$ ($r\ge0$, $j\ge1$) be a maximal run of zeros
of $M$ at consecutive integers, and let $\nu_-$ and $\nu_+$ be the non-null
excursions ending at $z_j$ and starting at $z_{j+r}$. The run is a
\emph{crossing} if $\sigma_{\nu_-}\sigma_{\nu_+}=-1$ and a \emph{touching}
run if $\sigma_{\nu_-}\sigma_{\nu_+}=+1$. At a touching run $M$ vanishes
without changing sign.
\end{definition}

\begin{definition}[Regularized skeleton]
\label{def:Ms}
The \emph{skeleton} $\Ms$ is the continuous, piecewise linear function
defined on each non-null excursion by
\begin{equation}
\label{eq:Ms_def}
  \Ms(t)=
  \begin{cases}
    \sigma_\nu\sqrt{H_\nu}\,\dfrac{t-u_\nu}{a_\nu}, & u_\nu\le t\le \tau_\nu,\\[2ex]
    \sigma_\nu\sqrt{H_\nu}\,\dfrac{v_\nu-t}{b_\nu}, & \tau_\nu\le t\le v_\nu ,
  \end{cases}
\end{equation}
and by $\Ms\equiv0$ on null excursions. On the terminal block $(u,N]$ of a
finite computation, if $M(N)\neq0$, the descent ends at
$\sigma\sqrt{|M(N)|}$ instead of at $0$. The \emph{unregularized skeleton}
$\Msk$ is defined in the same way with peak $\sigma_\nu H_\nu$ instead of
$\sigma_\nu\sqrt{H_\nu}$.
\end{definition}

\begin{definition}[The three curves]
\label{def:areas}
With $\Ms^{\pm}=\max(\pm \Ms,0)$,
\begin{equation}
\label{eq:areas_def}
  \Ap(x)=\int_0^x \Ms^{+}(t)\,dt,\qquad
  \Am(x)=-\int_0^x \Ms^{-}(t)\,dt,\qquad
  \Anet(x)=\Ap(x)+\Am(x)=\int_0^x \Ms(t)\,dt .
\end{equation}
In the figure of version 5, $\Ap$ is the blue curve, $\Am$ the red one and
$\Anet$ the black one. $\Ap$ is nondecreasing, $\Am$ is nonincreasing, and
each is constant on the excursions of the opposite sign and on null
excursions; this produces the plateaus.
\end{definition}

\section{Exact formulas}
\label{sec:formulas}

\noindent\emph{Status.} Original observations with complete proofs.
Lemma~\ref{lem:critical} corrects the corresponding lemma of version 5.

\begin{lemma}[Increment per excursion]
\label{lem:jump}
\emph{[unconditional]} For every complete excursion,
\begin{equation}
\label{eq:jump}
  \Anet(v_\nu)-\Anet(u_\nu)=\sigma_\nu\,\frac{\sqrt{H_\nu}\,\ell_\nu}{2}.
\end{equation}
The discrete version agrees:
\begin{equation}
\label{eq:jump_discrete}
  \sum_{n\in I_\nu}\Ms(n)
  =\sigma_\nu\sqrt{H_\nu}\,\frac{(a_\nu+1)+(b_\nu-1)}{2}
  =\sigma_\nu\sqrt{H_\nu}\,\frac{\ell_\nu}{2}.
\end{equation}
\end{lemma}

\begin{proof}
For null excursions both sides vanish. Otherwise the two pieces of $\Ms$ are
triangles of height $\sqrt{H_\nu}$ with bases $a_\nu$ and $b_\nu$. In the
discrete sum the ascent contributes
$\sqrt{H_\nu}\sum_{j=1}^{a_\nu}j/a_\nu=\sqrt{H_\nu}(a_\nu+1)/2$ and the
descent $\sqrt{H_\nu}\sum_{j=0}^{b_\nu-1}j/b_\nu=\sqrt{H_\nu}(b_\nu-1)/2$.
\end{proof}

\begin{proposition}[Explicit form of the three curves]
\label{prop:explicit}
\emph{[unconditional]} Let $x\in I_\nu$ with $I_\nu$ non-null. Then
\begin{equation}
\label{eq:explicit_A}
  \Anet(x)=
  \begin{cases}
    \Anet(u_\nu)+\sigma_\nu\sqrt{H_\nu}\dfrac{(x-u_\nu)^2}{2a_\nu},
      & u_\nu\le x\le\tau_\nu,\\[2.4ex]
    \Anet(u_\nu)+\sigma_\nu\sqrt{H_\nu}
      \left(\dfrac{a_\nu}{2}+(x-\tau_\nu)-\dfrac{(x-\tau_\nu)^2}{2b_\nu}\right),
      & \tau_\nu\le x\le v_\nu,
  \end{cases}
\end{equation}
and, at the zeros of $M$,
\begin{align}
  \Anet(v_k)&=\sum_{\nu\le k}\sigma_\nu\frac{\sqrt{H_\nu}\,\ell_\nu}{2},
  \label{eq:explicit_net}\\
  \Ap(v_k)&=\sum_{\substack{\nu\le k\\ \sigma_\nu=+1}}\frac{\sqrt{H_\nu}\,\ell_\nu}{2},
  \label{eq:explicit_plus}\\
  \Am(v_k)&=-\sum_{\substack{\nu\le k\\ \sigma_\nu=-1}}\frac{\sqrt{H_\nu}\,\ell_\nu}{2}.
  \label{eq:explicit_minus}
\end{align}
The three curves are piecewise quadratic, of class $C^1$, and determined by
the data $(u_\nu,\tau_\nu,v_\nu,H_\nu,\sigma_\nu)$.
\end{proposition}

\begin{proof}
Integrate \eqref{eq:Ms_def}; for the descent use $v_\nu-t=b_\nu-(t-\tau_\nu)$.
The formulas at the zeros follow from Lemma~\ref{lem:jump}. The derivative
$\Anet'=\Ms$ is continuous.
\end{proof}

\begin{lemma}[Critical points and inflections]
\label{lem:critical}
\emph{[unconditional]}
\begin{enumerate}[label=\textup{(\roman*)}]
\item $\Anet'=\Ms$; $\Anet$ is strictly monotone on each non-null excursion
and constant on each null excursion.
\item $\Anet$ has a local extremum at a run of zeros of $M$ if and only if the
run is a crossing (Definition~\ref{def:crossing_touching}). At a touching run
$\Anet$ is monotone through the run.
\item $\Anet''$ is piecewise constant. It changes sign at $\tau_\nu$ for every
non-null excursion and at every touching run, and nowhere else.
\end{enumerate}
\end{lemma}

\begin{proof}
(i) and (ii) follow from $\Anet'=\Ms$ and the sign of $\Ms$ on each
excursion. For (iii), on the ascent of $I_\nu$ one has
$\Anet''=\sigma_\nu\sqrt{H_\nu}/a_\nu$ and on the descent
$\Anet''=-\sigma_\nu\sqrt{H_\nu}/b_\nu$, which have opposite signs. Across a
run of zeros $\Anet''$ passes from $-\sigma_{\nu_-}\sqrt{H_{\nu_-}}/b_{\nu_-}$
to $\sigma_{\nu_+}\sqrt{H_{\nu_+}}/a_{\nu_+}$ (with value $0$ on intermediate
null excursions); the sign changes if and only if
$\sigma_{\nu_+}=\sigma_{\nu_-}$.
\end{proof}

\begin{remark}[Status of the version 5 statement]
\label{rem:critical_v5}
Version 5 stated that the local extrema of $\Anet$ are exactly the zeros of
$M$ and that the inflection points are exactly the first argmaxima $\tau_\nu$.
Both statements are \emph{incorrect} in the presence of touching runs: a
touching zero is a critical point that is not an extremum, and it is an
inflection point.
\end{remark}

\section{Internal correlations of an excursion}
\label{sec:internal_corr}

\noindent\emph{Status.} Elementary identities, original to this notebook,
with complete proofs. The inequality \eqref{eq:Q_upper} is new in version 6.

\medskip

The excursion decomposition reorganizes the quadratic form of $M$. Consider
first an arbitrary finite block $x_1,\dots,x_L\in\{-1,0,1\}$ with partial sums
and internal correlations
\begin{equation}
\label{eq:block_def}
  S_t=\sum_{j=1}^t x_j\quad(1\le t\le L),\qquad
  C_h=\sum_{j=1}^{L-h}x_jx_{j+h}\quad(1\le h<L).
\end{equation}

\begin{lemma}[Triangular correlation identity]
\label{lem:triangular-corr}
\emph{[unconditional]} For every finite block,
\begin{equation}
\label{eq:general-triangular}
  \sum_{t=1}^{L}S_t^2
  =S_L^2+S_L\sum_{t=1}^{L-1}S_t-\sum_{h=1}^{L-1}hC_h.
\end{equation}
In particular, for a bridge $S_L=0$,
\begin{equation}
\label{eq:bridge-triangular}
  \sum_{t=1}^{L-1}S_t^2=-\sum_{h=1}^{L-1}hC_h.
\end{equation}
\end{lemma}

\begin{proof}
One has
\begin{equation}
\label{eq:proof_triangular_max}
  \sum_{t=1}^{L}S_t^2
  =\sum_{i,j=1}^{L}x_ix_j\bigl(L+1-\max(i,j)\bigr).
\end{equation}
With $\max(i,j)=\tfrac12(i+j+|i-j|)$, the $|i-j|$ term gives
$-\sum_{h<L}hC_h$, and the $i+j$ term gives $-S_L\sum_i ix_i$, which is
reduced by
\begin{equation}
\label{eq:proof_triangular_abel}
  \sum_{i=1}^{L} i x_i
  =LS_L-\sum_{t=1}^{L-1}S_t.
\end{equation}
Collecting the terms $(L+1)S_L^2-L S_L^2$ yields
\eqref{eq:general-triangular}.
\end{proof}

For a complete non-null excursion $I_\nu=(u_\nu,v_\nu]$ put $L=\ell_\nu$ and
orient its increments by
\begin{equation}
\label{eq:orientation}
  x_j=\sigma_\nu\mu(u_\nu+j),\qquad
  S_t=\sum_{j=1}^{t}x_j=|M(u_\nu+t)|.
\end{equation}
Then $S_t>0$ for $1\le t<\ell_\nu$ and $S_{\ell_\nu}=0$. Write
\begin{equation}
\label{eq:Q_def}
  Q_\nu:=\sum_{t=1}^{\ell_\nu-1}S_t^2,
  \qquad
  C_h^{(\nu)}:=\sum_{j=1}^{\ell_\nu-h}x_jx_{j+h},
\end{equation}
with $Q_\nu=0$ for null excursions. Equation \eqref{eq:bridge-triangular}
gives the exact local identity
\begin{equation}
\label{eq:Qcorr}
  Q_\nu=-\sum_{h=1}^{\ell_\nu-1}hC_h^{(\nu)}\ge0.
\end{equation}

\begin{lemma}[Zeroth moment of the internal correlations]
\label{lem:zeroth-corr}
\emph{[unconditional]} If $q_\nu=\sum_{j=1}^{\ell_\nu}x_j^2$ is the number
of non-zero steps in the excursion, then
\begin{equation}
\label{eq:zeroth-corr}
  \sum_{h=1}^{\ell_\nu-1}C_h^{(\nu)}=-\frac{q_\nu}{2}.
\end{equation}
\end{lemma}

\begin{proof}
Since $\sum_jx_j=0$,
\begin{equation}
\label{eq:proof_zeroth}
  0=\Bigl(\sum_jx_j\Bigr)^2
   =\sum_jx_j^2+2\sum_{h=1}^{\ell_\nu-1}C_h^{(\nu)}.
\end{equation}
\end{proof}

\begin{proposition}[Quadratic decomposition by excursions]
\label{prop:quadratic-decomp}
\emph{[unconditional]} If $x$ is a zero of $M$, then
\begin{equation}
\label{eq:quadratic-decomp-zero}
  \sum_{n\le x}M(n)^2
   =\sum_{v_\nu\le x}Q_\nu
   =-\sum_{v_\nu\le x}\ \sum_{h=1}^{\ell_\nu-1}hC_h^{(\nu)}.
\end{equation}
For general $x$, with $u$ the last zero of $M$ not exceeding $x$,
\begin{equation}
\label{eq:quadratic-decomp-general}
  \sum_{n\le x}M(n)^2=\sum_{v_\nu\le x}Q_\nu+Q_{\rm term}(x),
  \qquad
  Q_{\rm term}(x):=\sum_{u<n\le x}M(n)^2,
\end{equation}
and $Q_{\rm term}(x)$ is given by \eqref{eq:general-triangular} applied to
the terminal block.
\end{proposition}

\begin{remark}[Blindness to pairs crossing excursion boundaries]
\label{rem:blindness-quadratic}
Globally, for integer $x$,
\begin{equation}
\label{eq:global_quadratic}
  \sum_{n\le x}M(n)^2
  =\sum_{a,b\le x}\mu(a)\mu(b)\bigl(x+1-\max(a,b)\bigr),
\end{equation}
which contains all pairs $a,b\le x$. For $a\in I_\nu$ and $b\in I_\kappa$
with $\nu<\kappa$ the weight depends only on $b$, and
$\sum_{a\in I_\nu}\mu(a)=0$; hence all pairs crossing a zero cancel exactly,
and Proposition~\ref{prop:quadratic-decomp} retains only pairs inside one
excursion (plus the terminal block).
\end{remark}

Let
\begin{equation}
\label{eq:beta_def}
  \beta_\nu:=\sum_{t=1}^{\ell_\nu-1}S_t
             =\sum_{n\in I_\nu}|M(n)|.
\end{equation}
By Cauchy--Schwarz,
\begin{equation}
\label{eq:beta-Q}
  \beta_\nu^2\le \ell_\nu Q_\nu
  =-\ell_\nu\sum_{h=1}^{\ell_\nu-1}hC_h^{(\nu)}.
\end{equation}
Thus the first weighted moment of the internal correlations controls the area
of the excursion. The converse needs information on the flatness of the
profile.

To separate the support from the sign geometry, let
$P_h^{(\nu)}:=\#\{1\le j\le \ell_\nu-h:\ x_jx_{j+h}\ne0\}$ and let
$D_h^{(\nu)}$ be the number of those pairs with opposite signs. Then
\begin{equation}
\label{eq:C-P-D}
  C_h^{(\nu)}=P_h^{(\nu)}-2D_h^{(\nu)},
\end{equation}
and hence
\begin{equation}
\label{eq:Q-P-D}
  Q_\nu=\sum_{h=1}^{\ell_\nu-1}
  h\bigl(2D_h^{(\nu)}-P_h^{(\nu)}\bigr).
\end{equation}

\begin{remark}[Height versus quadratic energy]
\label{rem:height-Q}
\emph{[unconditional]} If a non-null excursion reaches height $H_\nu$, the
least possible value of $Q_\nu$ is obtained by climbing and descending at
maximal speed with no zero steps, so that
\begin{equation}
\label{eq:height_Q_min}
  Q_\nu\ge 2\sum_{r=1}^{H_\nu-1}r^2+H_\nu^2
  =\frac{2H_\nu^3+H_\nu}{3}.
\end{equation}
In the other direction, $S_t\le\min(t,\ell_\nu-t)\le\ell_\nu/2$ gives
\begin{equation}
\label{eq:Q_upper}
  Q_\nu\le\frac{(\ell_\nu-1)\ell_\nu^2}{4}\le\frac{\ell_\nu^3}{4}.
\end{equation}
A local bound $Q_\nu\ll \ell_\nu^{2+o(1)}$ therefore gives only
$H_\nu\ll \ell_\nu^{2/3+o(1)}$, and the natural primary object is the area
normalization $\beta_\nu/\ell_\nu^{3/2}$ rather than the height
normalization $H_\nu/\sqrt{\ell_\nu}$.
\end{remark}

\section{Where \texorpdfstring{$\zeta$}{zeta} enters}
\label{sec:zeta}

\noindent\emph{Status.} \eqref{eq:explicitM} is standard
\cite[\S14.27]{Titchmarsh1986}. \eqref{eq:explicitA} is obtained formally
by the same contour method; its residues are checked, the contour details
are not written out. Proposition~\ref{prop:dipoles} is original, with proof.

\medskip

The black curve depends on $M$ nonlinearly, through $\sqrt{H_\nu}$ and
through the selection of $u_\nu,\tau_\nu,v_\nu$. No closed expression of its
Mellin transform in terms of $\zeta$ is known to the author (version 5
asserted that none exists; that assertion is unproved and is withdrawn). The
nodes $u_\nu,v_\nu$ are zeros of $M$, and \emph{[under RH and simple zeros]}
\begin{equation}
\label{eq:explicitM}
  M(x)=\sum_{\rho}\frac{x^\rho}{\rho\,\zeta'(\rho)}-2
  +\sum_{n\ge1}\frac{(-1)^{n-1}(2\pi/x)^{2n}}{(2n)!\,n\,\zeta(2n+1)},
\end{equation}
for non-integer $x$, the sum over $\rho$ taken along a suitable sequence
$|\gamma|<T_k\to\infty$ \cite[\S14.27]{Titchmarsh1986}. For
$\widetilde{A}(x):=\int_0^xM(t)\,dt$ one has, for $\operatorname{Re}s>1$,
\begin{equation}
\label{eq:mellin_A}
  \int_1^\infty \widetilde{A}(x)x^{-s-2}\,dx=\frac{1}{s(s+1)\zeta(s)}.
\end{equation}
Shifting the contour as in \cite[\S14.27]{Titchmarsh1986} and collecting the
residues at $\rho$, at $s=0$ ($\zeta(0)=-\tfrac12$) and at $s=-1$
($\zeta(-1)=-\tfrac1{12}$), one formally arrives at \emph{[under RH and
simple zeros]}
\begin{equation}
\label{eq:explicitA}
  \widetilde{A}(x)=\sum_\rho\frac{x^{\rho+1}}{\rho(\rho+1)\zeta'(\rho)}
  -2x+12+O(x^{-1}).
\end{equation}
The factor $(\rho+1)^{-1}$ improves convergence subject to the size of
$1/\zeta'(\rho)$; hypotheses of Gonek--Hejhal type controlling this are
discussed in \cite{Ng2004}. The area of $\Msk$ differs from
$\widetilde{A}$ by explicit local terms:

\begin{proposition}[Dipole identity]
\label{prop:dipoles}
\emph{[unconditional]} For each excursion let its \emph{legs} be the integer
sets $\lambda=(u_\nu,\tau_\nu]\cap\mathbb{Z}$ and $(\tau_\nu,v_\nu]\cap\mathbb{Z}$,
with midpoints $c_\lambda=(u_\nu+\tau_\nu)/2$ and $(\tau_\nu+v_\nu)/2$. Then
for every zero $v_k$,
\begin{equation}
\label{eq:dipoles}
  \int_0^{v_k}\Msk(t)\,dt=\widetilde{A}(v_k)
  +\sum_{\lambda\subset(0,v_k]}\ \sum_{m\in\lambda}\mu(m)\,(m-c_\lambda).
\end{equation}
\end{proposition}

\begin{proof}
On an ascent leg, $\int_{u}^{\tau}\Msk=\sigma H a/2$ and
$\int_u^\tau M=\sum_{m=u+1}^{\tau}\mu(m)(\tau-m)$, while
$\sigma H=\sum_{m=u+1}^{\tau}\mu(m)$; the difference is
$\sum_m\mu(m)(a/2-\tau+m)=\sum_m\mu(m)(m-c_\lambda)$. On a descent leg,
$\int_\tau^v\Msk=\sigma Hb/2$ and
$\int_\tau^vM=b\,\sigma H+\sum_{m=\tau+1}^{v}\mu(m)(v-m)$ with
$\sum_{m=\tau+1}^{v}\mu(m)=-\sigma H$; the difference is
$\sum_m\mu(m)(m-(v-b/2))$. Sum over legs.
\end{proof}

\section{Bounds}
\label{sec:bounds}

\noindent\emph{Status.} Original, elementary; the input bounds for $M$ are
standard \cite{Walfisz1963,OdlyzkoTeRiele1985,KotnikTeRiele2006}.
Lemma~\ref{lem:transfer} corrects the hypothesis of version 5;
Remarks~\ref{rem:onedirection} and~\ref{rem:omega} correct overstatements.

\begin{lemma}[Transfer]
\label{lem:transfer}
\emph{[unconditional]} Let $x\ge1$ and let $x^*\ge x$ be the first zero of
$M$ not smaller than $x$ (or $N$ in the terminal block). If
$|M(n)|\le c\,n^{1/2}$ for all $1\le n\le x^*$ (so necessarily $c\ge1$),
then
\begin{equation}
\label{eq:transfer_bound}
  |\Ms(t)|\le c^{1/2}t^{1/4}\quad(0\le t\le x),\qquad
  \Ap(x)+|\Am(x)|\le\tfrac45\,c^{1/2}x^{5/4}.
\end{equation}
\end{lemma}

\begin{proof}
For $t\ge\tau_\nu$, $|\Ms(t)|\le\sqrt{H_\nu}\le c^{1/2}\tau_\nu^{1/4}\le
c^{1/2}t^{1/4}$. For $u_\nu\le t\le\tau_\nu$ the chord from $(u_\nu,0)$ to
$(\tau_\nu,\sqrt{H_\nu})$ lies below the concave function $c^{1/2}t^{1/4}$,
which is $\ge0$ at $u_\nu$ and $\ge\sqrt{H_\nu}$ at $\tau_\nu\le x^*$.
Integrate.
\end{proof}

\begin{corollary}
\label{cor:uncond}
\emph{[unconditional]} From the Vinogradov--Korobov bound
\cite{Walfisz1963}
\begin{equation}
\label{eq:VK_M}
  M(x)\ll x\exp\bigl(-c(\log x)^{3/5}(\log\log x)^{-1/5}\bigr)
\end{equation}
one obtains
\begin{equation}
\label{eq:uncond_area}
  \Ap(x)+|\Am(x)|\ll x^{3/2}\exp\bigl(-\tfrac c2(\log x)^{3/5}(\log\log x)^{-1/5}\bigr).
\end{equation}
\end{corollary}

\begin{corollary}
\label{cor:rh}
\emph{[unconditional implication]} If RH holds, then
\begin{equation}
\label{eq:rh_area}
  \Ap(x)+|\Am(x)|\ll x^{5/4+\eps}.
\end{equation}
Consequently, if $\Ap(x)+|\Am(x)|\gg x^{5/4+\eps}$ along a sequence, RH is
false.
\end{corollary}

\begin{remark}[One direction only, as far as known]
\label{rem:onedirection}
Since $\ell_\nu\ge2H_\nu$, a bound $\sqrt{H_\nu}\ell_\nu\ll x^{5/4+\eps}$
yields only $H_\nu\ll x^{5/6+\eps}$, hence only the zero-free half-plane
$\operatorname{Re}s>5/6$. Version 5 called the black curve ``a refutation
test, not a criterion''. Whether \eqref{eq:rh_area} is equivalent to RH is
\emph{unresolved}; the direct route fails, which is all that is shown.
\end{remark}

\begin{remark}[Oscillation of the heights]
\label{rem:omega}
Since $\limsup M(x)x^{-1/2}>1.06$ and $\liminf M(x)x^{-1/2}<-1.009$
\cite{OdlyzkoTeRiele1985} (improved in \cite{KotnikTeRiele2006}), there are
infinitely many $n$ with $M(n)>1.06\sqrt n$; the excursion containing $n$
has $u_\nu<n$, so $H_\nu>1.06\,u_\nu^{1/2}$ for infinitely many positive
excursions, and similarly for negative ones. Version 5 claimed
$H_\nu\gg v_\nu^{1/2-\eps}$; this is \emph{not justified}, since $v_\nu$
may be much larger than $n$.
\end{remark}

\section{Criteria equivalent to RH}
\label{sec:criteria}

\noindent\emph{Status.} The Mellin identity \eqref{eq:mellin_M} is standard
\cite{MontgomeryVaughan2007}; the equivalences, pruning results and
Corollary~\ref{cor:pruned_local} are original, with proofs.
Heuristics~\ref{heu:renewal} and~\ref{heu:local_fails} are heuristic.

\medskip

For $\operatorname{Re}s>1$,
\begin{equation}
\label{eq:mellin_M}
  \int_1^\infty M(x)x^{-s-1}\,dx=\frac{1}{s\,\zeta(s)} .
\end{equation}
If $\int_1^\infty|M(x)|x^{-\sigma-1}dx<\infty$ for every $\sigma>1/2$, the
left side continues analytically to $\operatorname{Re}s>1/2$, so $\zeta$ has
no zeros there.

\begin{proposition}[$L^1$ criterion]
\label{prop:L1}
\emph{[unconditional]} RH holds if and only if, for every $\eps>0$,
\begin{equation}
\label{eq:L1_criterion}
  \sum_{v_\nu\le x}\beta_\nu+\beta_{\rm term}(x)=\sum_{n\le x}|M(n)|
  \ll x^{3/2+\eps}.
\end{equation}
\end{proposition}

\begin{proof}
RH gives $M(x)\ll x^{1/2+\eps}$ \cite[\S14.25]{Titchmarsh1986}. Conversely,
by partial summation \eqref{eq:L1_criterion} gives
$\sum_n|M(n)|n^{-\sigma-1}<\infty$ for $\sigma>1/2+\eps$, and $M$ is
constant on $[n,n+1)$; apply the remark after \eqref{eq:mellin_M}.
\end{proof}

\begin{lemma}[Pruning]
\label{lem:pruning}
\emph{[unconditional]} $\beta_\nu\le H_\nu\ell_\nu\le\ell_\nu^2/2$, hence
\begin{equation}
\label{eq:pruning}
  \sum_{v_\nu\le x,\ \ell_\nu\le L_0}\beta_\nu\le\frac{xL_0}{2}.
\end{equation}
\end{lemma}

\begin{corollary}[Reduction to long excursions]
\label{cor:long}
\emph{[unconditional]} For fixed $\delta\ge0$ and $L_0=x^{1/2-\delta}$, RH
holds if and only if
\begin{equation}
\label{eq:long_reduction}
  \sum_{v_\nu\le x,\ \ell_\nu>L_0}\beta_\nu+\beta_{\rm term}(x)\ll x^{3/2+\eps}
  \quad\text{for every }\eps>0,
\end{equation}
where the terminal block is included when its length exceeds $L_0$. At most
$x/L_0$ excursions are involved.
\end{corollary}

\begin{remark}[The $\sqrt{x}$ frontier]
\label{rem:sqrt_frontier}
The cut may even depend on $\eps$: with $L_0=x^{1/2+\eps/2}$ the discarded
part is $\le x^{3/2+\eps/2}/2$. The frontier is therefore $x^{1/2+o(1)}$;
what requires new information is a cut $x^{\theta}$ with a fixed
$\theta>1/2$ (version 5 stated the frontier as $x^{1/2}$).
\end{remark}

\begin{heuristic}[Renewal tail]
\label{heu:renewal}
This is a heuristic argument. If excursion lengths behave like those of a
recurrent random walk,
\begin{equation}
\label{eq:renewal_tail}
  \#\{v_\nu\le x:\ \ell_\nu>L\}\asymp\sqrt{x}\,L^{-1/2},
\end{equation}
and with $\beta_\nu\asymp c_\beta\ell_\nu^{3/2}$ the dyadic range
$[L,2L]$ contributes $\asymp\sqrt{x}\,L$ to \eqref{eq:L1_criterion}, so the
terminal scales dominate.
\end{heuristic}

\begin{remark}[Limit of pruning]
\label{rem:pruning_limit}
For $\ell_\nu\asymp x$ the trivial bound gives $\beta_\nu\ll x^2$ against a
budget $x^{3/2+\eps}$: a factor $x^{1/2}$ is missing, and there are only
$O(1)$ such excursions.
\end{remark}

\begin{theorem}[Correlation criterion]
\label{thm:corr-criterion}
\emph{[unconditional]} RH holds if and only if, for every $\eps>0$,
\begin{equation}
\label{eq:RH-corr}
  \sum_{n\le x}M(n)^2=\sum_{v_\nu\le x}Q_\nu+Q_{\rm term}(x)
  =-\sum_{v_\nu\le x}\sum_{h}hC^{(\nu)}_h+Q_{\rm term}(x)\ll x^{2+\eps}.
\end{equation}
\end{theorem}

\begin{proof}
RH gives $M\ll x^{1/2+\eps}$. Conversely Cauchy--Schwarz gives
$\sum_{n\le x}|M(n)|\le x^{1/2}(\sum M^2)^{1/2}\ll x^{3/2+\eps/2}$; apply
Proposition~\ref{prop:L1}.
\end{proof}

\begin{corollary}[Local sufficient condition, version 5]
\label{cor:local-corr-sufficient}
\emph{[unconditional implication]} If for every $\eps>0$ there is
$C_\eps$ with
\begin{equation}
\label{eq:local_Q}
  Q_\nu\le C_\eps\,\ell_\nu^{2+\eps}\quad\text{for every excursion
  (including the terminal block)},
\end{equation}
then RH holds.
\end{corollary}

\begin{heuristic}[The hypothesis of Corollary~\ref{cor:local-corr-sufficient} fails]
\label{heu:local_fails}
This is a heuristic argument. In a model where, after a zero of $M$, the next
values of $\mu$ behave locally at random with squarefree density $6/\pi^2$,
the pattern ``$k$ non-zero steps of one sign followed by $k$ of the other''
follows a given zero with probability $\ge c_0^{-k}$. Since $M$ has
infinitely many zeros, excursions with $H_\nu=k$ and $\ell_\nu\asymp k$
should occur with $k\to\infty$. By \eqref{eq:height_Q_min}
$Q_\nu/\ell_\nu^{2+\eps}\gg k^{1-\eps}\to\infty$, so \eqref{eq:local_Q}
fails for every $\eps\in(0,1)$; for $\eps\ge1$ it is trivially true by
\eqref{eq:Q_upper}. Version 5 called the condition ``sufficient but not
weak''; heuristically it is never satisfied.
\end{heuristic}

\begin{corollary}[Pruned local condition]
\label{cor:pruned_local}
\emph{[unconditional implication]} If for every $\eps>0$ there is $C_\eps$
with
\begin{equation}
\label{eq:pruned_Q}
  Q_\nu\le C_\eps\,\ell_\nu^{2+\eps}\quad\text{whenever }\ell_\nu>v_\nu^{1/2}
\end{equation}
(and likewise for the terminal block when its length exceeds $x^{1/2}$),
then RH holds.
\end{corollary}

\begin{proof}
By \eqref{eq:Q_upper}, excursions with $\ell_\nu\le v_\nu^{1/2}\le x^{1/2}$
satisfy $Q_\nu\le\ell_\nu^3/4\le\ell_\nu x/4$, and these sum to at most
$x^2/4$. The others contribute $\le C_\eps\sum\ell_\nu^{2+\eps}\le
C_\eps x^{2+\eps}$. Apply Theorem~\ref{thm:corr-criterion}.
\end{proof}

\noindent The ballistic excursions of Heuristic~\ref{heu:local_fails} have
$\ell_\nu\asymp k\ll v_\nu^{1/2}$ and do not contradict \eqref{eq:pruned_Q}.

\begin{remark}[Scale of the correlation sums]
\label{rem:chowla-scale}
Under the crude model in which ``$j$ squarefree'' and ``$j+h$ squarefree''
are independent with probability $6/\pi^2$ each,
$P_h\approx(6/\pi^2)^2(L-h)$ and
\begin{equation}
\label{eq:chowla_crude}
  \sum_{h<L}hP_h\approx\Bigl(\frac{6}{\pi^2}\Bigr)^2\sum_{h<L}h(L-h)
  \approx\frac{6}{\pi^4}L^3\approx0.0616\,L^3 .
\end{equation}
Version 5 gave $\tfrac12(6/\pi^2)^2L^3\approx0.18L^3$, omitting the factor
$L-h$; that constant is \emph{incorrect}, the order $L^3$ stands. For
$|\sum_hhC_h|$ of size $B(\sum_hh^2)^{1/2}\asymp BL^{3/2}$ to be $O(L^2)$ one
needs $B=O(\sqrt L)$ (a calibration, not a proof).
\end{remark}

\begin{remark}[Conditioning on the bridge]
\label{rem:conditioning}
Inside an excursion the correlations are constrained by
\eqref{eq:Qcorr}, \eqref{eq:zeroth-corr} and by the geometry of the bridge;
with the defect $d_\nu:=\ell_\nu/2-H_\nu\ge0$, a third constraint of this
type is a research direction, not a result.
\end{remark}

\begin{theorem}[Energy criterion]
\label{thm:energy}
\emph{[unconditional]} Define, with $u$ the last zero $\le x$ and
$H_{\rm run}(x)=\max_{u<n\le x}|M(n)|$,
\begin{equation}
\label{eq:energy_def}
  \En(x):=\frac14\Bigl(\sum_{v_\nu\le x}H_\nu\ell_\nu+(x-u)H_{\rm run}(x)\Bigr).
\end{equation}
For a complete excursion $\tfrac14H_\nu\ell_\nu=\int_{I_\nu}(J')^2$, where
$J$ interpolates $\Anet$ linearly between consecutive zeros; for the
terminal block this interpretation does not hold and \eqref{eq:energy_def} is
a definition. Then
\begin{equation}
\label{eq:energy_sandwich}
  \sum_{n\le x}|M(n)|\le4\En(x)\le x\max_{n\le x}|M(n)|,
\end{equation}
and RH holds if and only if $\En(x)\ll x^{3/2+\eps}$ for every $\eps>0$.
\end{theorem}

\begin{proof}
$\beta_\nu\le H_\nu\ell_\nu$ and $\sum\ell_\nu\le x$; apply
Proposition~\ref{prop:L1}.
\end{proof}

\section{Trial-division horizon of an excursion}
\label{sec:prime_horizon}

\noindent\emph{Status.} Lemma~\ref{lem:square-horizon} is standard (trial
division). Warning~\ref{warn:information_language},
Lemma~\ref{lem:rough_part} and Proposition~\ref{prop:double-exp} are
original and elementary. Question~\ref{q:finiteness_expanding} is open.

\begin{lemma}[Square horizon]
\label{lem:square-horizon}
\emph{[unconditional]} Let $u\ge2$. Trial division by the primes $p\le u$
alone determines $\mu(n)$ for every $n\le u^2$.
\end{lemma}

\begin{proof}
Remove the prime factors $p\le u$ of $n$; if one occurs twice, $\mu(n)=0$.
The cofactor $q\le n\le u^2$ has all prime factors $>u$, so it is $1$ or a
single prime (two such factors would give $q>u^2$).
\end{proof}

\begin{warning}[On the word ``information'']
\label{warn:information_language}
Version 5 described primes $>u$ as information not available at the start
$u$ of an excursion. In the logical sense this is \emph{incorrect}: the
primes $\le u$ determine, by the sieve of Eratosthenes, all primes $\le u^2$,
and by iteration all primes. What Lemma~\ref{lem:square-horizon} measures is
\emph{sieve depth}. In Sections~\ref{sec:prime_horizon}
to~\ref{sec:local_large_prime} the words ``known'', ``new'' and
``information'' are used only in this sense.
\end{warning}

\begin{definition}[Smooth and rough parts]
\label{def:smooth_rough}
For $y\ge1$ write $n=s_y(n)\,r_y(n)$, where every prime factor of $s_y(n)$
is $\le y$ and every prime factor of $r_y(n)$ is $>y$. Let $\Pp(n)$ be the
largest prime factor of $n$, with $\Pp(1)=1$.
\end{definition}

\begin{lemma}
\label{lem:rough_part}
\emph{[unconditional]}
\begin{equation}
\label{eq:mu_split}
  \mu(n)=\mu\bigl(s_y(n)\bigr)\,\mu\bigl(r_y(n)\bigr);
\end{equation}
if $n<y^k$ then $\Omega(r_y(n))\le k-1$, and if $n\le y^2$ then $r_y(n)$ is
$1$ or a prime.
\end{lemma}

\begin{proof}
The two parts are coprime. If $\Omega(r)\ge1$ then $y^{\Omega(r)}<r\le n$.
\end{proof}

\begin{definition}[Self-contained and expanding excursions]
\label{def:info-excursion}
An excursion is \emph{information-self-contained} if $v_\nu\le u_\nu^2$ and
\emph{information-expanding} if $v_\nu>u_\nu^2$ (names kept from version 5,
read in the sense of Warning~\ref{warn:information_language}). The fraction
of $I_\nu$ below $u_\nu^2$ is
\begin{equation}
\label{eq:known-fraction}
  \vartheta_\nu=\min\Bigl\{1,\frac{u_\nu^2-u_\nu}{\ell_\nu}\Bigr\}.
\end{equation}
\end{definition}

\begin{proposition}[Double-exponential separation]
\label{prop:double-exp}
\emph{[unconditional]} Let $a_1<a_2<\dots$ be the starting points $\ge2$ of
information-expanding excursions. Then
\begin{equation}
\label{eq:double-exp-separation}
  a_{j+1}>a_j^2,\qquad a_j\ge a_1^{2^{j-1}}\ (j\ge1),
\end{equation}
with strict inequality in the second relation for $j\ge2$ (version 5 stated
it strictly for all $j$, which fails at $j=1$). Hence
\begin{equation}
\label{eq:loglog-exception-count}
  N_{\rm new}(X):=\#\{j:\ a_j\le X\}\le1+\log_2\frac{\log X}{\log2}.
\end{equation}
\end{proposition}

\begin{proof}
An expanding excursion starting at $a_j$ ends beyond $a_j^2$, and the next
one starts at or after that end.
\end{proof}

\begin{question}
\label{q:finiteness_expanding}
Are there only finitely many information-expanding excursions?
Equivalently: does $M$ vanish in $(x,x^2]$ for every sufficiently large $x$?
(If the zeros $u\ge x_0$ all have a zero in $(u,u^2]$, apply this to the
last zero $\le x$; the converse is immediate.) No unconditional result of
this form has been located; results counting sign changes of arithmetic
error terms and lower bounds for the mean of $|M|$ do not give it.
\end{question}

\begin{remark}[The horizon $u^2$ is not the multiplicative threshold of an excursion]
\label{rem:horizon_inactive}
For a self-contained excursion $\sqrt{v_\nu}\le u_\nu$. Every $n\le v_\nu$
has at most one prime factor $>\sqrt{v_\nu}$, so the threshold that governs
the multiplicative structure inside $I_\nu$ is $\sqrt{v_\nu}$, which lies
below $u_\nu$. The split ``known/new'' at $u_\nu$ does not match it, and for
a self-contained excursion the regime $t>u_\nu^2$ of
Section~\ref{sec:prime_copy_propagation} lies beyond $v_\nu$. That section
is therefore about $M(t)$, not about a single excursion;
Section~\ref{sec:local_large_prime} gives the local version.
\end{remark}

\section{Large-prime decomposition at level \texorpdfstring{$u^2$}{u^2}}
\label{sec:prime_copy_propagation}

\noindent\emph{Status.} All identities are original and elementary. The
asymptotic statements use only the prime number theorem with classical error
term and the Vinogradov--Korobov bound
\cite{MontgomeryVaughan2007,Walfisz1963}, and are proved here. Several
statements of version 5 are corrected; see
Warnings~\ref{warn:v5_Ny}, \ref{warn:v5_main_term},
\ref{warn:margin_vacuous} and~\ref{warn:v5_calibration}.

\begin{definition}[Baseline and large-prime part]
\label{def:baseline_channel}
For $u\ge2$ and $t\ge1$ put
\begin{equation}
\label{eq:baseline_channel_def}
  B_u(t):=\sum_{\substack{n\le t\\ \Pp(n)\le u^2}}\mu(n),\qquad
  E_u(t):=\sum_{\substack{n\le t\\ \Pp(n)>u^2}}\mu(n),\qquad
  M(t)=B_u(t)+E_u(t).
\end{equation}
Version 5 used an undefined ``baseline'' $B_u(t)$; \eqref{eq:baseline_channel_def}
is the definition adopted here. In the window $u^2<t<u^4$, $E_u$ is the
quantity $E_u^{(1)}$ of version 5.
\end{definition}

\begin{proposition}[Single-prime copy]
\label{prop:single_prime_copy}
\emph{[unconditional]} For a prime $p$ and $t<p^2$,
\begin{equation}
\label{eq:single-prime-copy}
  \Delta_p(t):=\sum_{a\le t/p}\mu(ap)=-M(t/p).
\end{equation}
Only $p>\sqrt t$ is used; the condition $p>u^2$ of version 5 plays no role.
\end{proposition}

\begin{proof}
$a\le t/p<p$, so $p\nmid a$ and $\mu(ap)=-\mu(a)$.
\end{proof}

\begin{proposition}[Sign content of one prime cone]
\label{prop:prime-cone-signs}
\emph{[unconditional]} Let $\Sq(y)=\sum_{a\le y}\mu(a)^2$. Among $ap$,
$a\le y<p$, the numbers of values $+1,-1,0$ of $\mu$ are
\begin{equation}
\label{eq:cone_counts}
  N_+=\frac{\Sq(y)-M(y)}{2},\qquad N_-=\frac{\Sq(y)+M(y)}{2},\qquad
  N_0=\lfloor y\rfloor-\Sq(y),
\end{equation}
with limiting proportions $3/\pi^2,\ 3/\pi^2,\ 1-6/\pi^2$.
\end{proposition}

\begin{proposition}[Linear window]
\label{prop:linear-new-primes}
\emph{[unconditional]} If $u^2<t<u^4$, no $n\le t$ has two prime factors
$>u^2$ (counted with multiplicity), and
\begin{equation}
\label{eq:first-order-total}
  E_u(t)=-\sum_{u^2<p\le t}M(t/p).
\end{equation}
\end{proposition}

\begin{proposition}[Prime-count form]
\label{prop:prime_count_form}
\emph{[unconditional]} For $u^2<t<u^4$ and $R:=t/u^2$,
\begin{equation}
\label{eq:prime_count_form}
  E_u(t)=-\sum_{a\le R}\mu(a)\bigl(\pi(t/a)-\pi(u^2)\bigr),
  \qquad
  M(t)=B_u(t)-\sum_{a\le R}\mu(a)\pi(t/a)+M(R)\pi(u^2).
\end{equation}
\end{proposition}

\begin{proof}
Expand \eqref{eq:first-order-total} as
$-\sum_{p}\sum_{a\le t/p}\mu(a)$ and interchange: for fixed $a$ the prime
runs over $u^2<p\le t/a$, which is non-empty only if $a<R$.
\end{proof}

\begin{remark}[How the primes above $u^2$ enter]
\label{rem:unknown_primes_enter}
The only input from primes $>u^2$ is the value of $\pi$ at the points $t/a$,
$a\le R$, weighted by $\mu(a)$ at the small scale $R<u^2$. Individual
positions of those primes matter only through these counts.
\end{remark}

\begin{definition}[Copy multiplicities]
\label{def:copy-multiplicity}
For $u^2<t<u^4$, $R=t/u^2$ and $1\le y\le\lfloor R\rfloor$ put
\begin{equation}
\label{eq:Ny-exact}
  N_y(t,u)=\#\Bigl\{p:\ \max\bigl(u^2,\tfrac{t}{y+1}\bigr)<p\le\tfrac ty\Bigr\},
  \qquad
  E_u(t)=-\sum_{1\le y\le R}N_y(t,u)M(y).
\end{equation}
\end{definition}

\begin{remark}[Asymptotics of the multiplicities, corrected]
\label{rem:copy-count-asymptotic}
For fixed $R>1$ and $u\to\infty$,
\begin{equation}
\label{eq:Ny-fixed-R}
  N_y\sim\frac{Ru^2}{2\log u}\cdot\frac1{y(y+1)}\quad(1\le y\le R-1),
  \qquad
  N_{\lfloor R\rfloor}\sim\frac{u^2}{2\log u}\Bigl(\frac R{\lfloor R\rfloor}-1\Bigr)
  \quad(R\notin\mathbb{Z}),
\end{equation}
and $N_R=0$ if $R\in\mathbb{Z}$.
\end{remark}

\begin{warning}
\label{warn:v5_Ny}
Version 5 stated the first relation of \eqref{eq:Ny-fixed-R} for all $y<R$.
This is \emph{incorrect} for $y=\lfloor R\rfloor$ when $R\notin\mathbb{Z}$,
since then $t/(y+1)<u^2$.
\end{warning}

\begin{definition}[Coherent coefficient]
\label{def:coherent_coefficient}
\begin{equation}
\label{eq:C_def}
  \Cc(R):=\int_1^R\frac{M(w)}{w^2}\,dw
  =\sum_{a\le R}\mu(a)\Bigl(\frac1a-\frac1R\Bigr)
  =\sum_{a\le R}\frac{\mu(a)}a-\frac{M(R)}R .
\end{equation}
\end{definition}

\begin{proposition}[Fixed $R$, corrected]
\label{prop:fixed_R_corrected}
\emph{[unconditional]} For fixed $R>1$, as $u\to\infty$,
\begin{equation}
\label{eq:E_fixed_R}
  E_u(Ru^2)=-\frac{Ru^2}{2\log u}\bigl(\Cc(R)+o(1)\bigr).
\end{equation}
\end{proposition}

\begin{proof}
In \eqref{eq:prime_count_form} the sum is finite; for $a<R$ the prime number
theorem gives $\pi(Ru^2/a)-\pi(u^2)\sim(R/a-1)u^2/(2\log u)$, and
$\sum_{a\le R}\mu(a)(R/a-1)=R\,\Cc(R)$.
\end{proof}

\begin{lemma}[The coefficient of version 5]
\label{lem:copy-weight-coefficient}
\emph{[unconditional]} The coefficient $A(R)$ of version 5, here
\begin{equation}
\label{eq:Theta_def}
  \Theta(R):=\sum_{1\le y\le R}\frac{M(y)}{y(y+1)}
  =\sum_{n\le R}\frac{\mu(n)}n-\frac{M(R)}{\lfloor R\rfloor+1},
\end{equation}
satisfies
\begin{equation}
\label{eq:Theta_C}
  \Theta(R)=\Cc\bigl(\lfloor R\rfloor+1\bigr),\qquad
  \Cc(R)=\Theta(R)-M(R)\Bigl(\frac1R-\frac1{\lfloor R\rfloor+1}\Bigr),
\end{equation}
and
\begin{equation}
\label{eq:C_tail}
  \Cc(R)=-\int_R^\infty\frac{M(w)}{w^2}\,dw,\qquad
  \Theta(R)=-\sum_{n>R}\frac{M(n)}{n(n+1)}.
\end{equation}
Version 5 wrote $M(R)/(R+1)$ in \eqref{eq:Theta_def}; this is correct only
for integer $R$.
\end{lemma}

\begin{proof}
$M$ is constant on $[y,y+1)$ and $\int_y^{y+1}w^{-2}dw=1/(y(y+1))$. By
\eqref{eq:mellin_M} and dominated convergence as $s\downarrow1$ (using
$M(w)\ll we^{-c\sqrt{\log w}}$), $\int_1^\infty M(w)w^{-2}dw=0$.
\end{proof}

\begin{warning}[Main term of version 5]
\label{warn:v5_main_term}
Version 5 asserted $E_u^{(1)}(Ru^2)\sim-\frac{Ru^2}{2\log u}\Theta(R)$. By
\eqref{eq:E_fixed_R} and \eqref{eq:Theta_C} this is \emph{correct only when
$M(\lfloor R\rfloor)=0$}. For $1<R<2$ one has
$E_u=-(\pi(t)-\pi(u^2))\sim-(R-1)\frac{u^2}{2\log u}$, whereas version 5
gives $-\frac R2\cdot\frac{u^2}{2\log u}$.
\end{warning}

\begin{corollary}[Decay of the coefficients]
\label{cor:AR-VK}
\emph{[unconditional]} By \eqref{eq:C_tail} and \eqref{eq:VK_M},
\begin{equation}
\label{eq:AR-VK}
  \Cc(R),\ \Theta(R)\ll\exp\bigl(-c_1(\log R)^{3/5}(\log\log R)^{-1/5}\bigr).
\end{equation}
\end{corollary}

\begin{proposition}[Both parts are large]
\label{prop:both_parts_large}
\emph{[unconditional]}
\begin{enumerate}[label=\textup{(\roman*)}]
\item For fixed $R>1$ with $\Cc(R)\ne0$,
\begin{equation}
\label{eq:E_fixed_R_size}
  E_u(Ru^2)\sim-\frac{Ru^2}{2\log u}\Cc(R),\qquad
  B_u(Ru^2)\sim+\frac{Ru^2}{2\log u}\Cc(R);
\end{equation}
for $1<R<2$, $\Cc(R)=1-1/R>0$.
\item There is $C>0$ such that, as $t\to\infty$, uniformly for $u^2<t<u^4$
with $R=t/u^2\ge\exp(C(\log\log t)^2)$,
\begin{equation}
\label{eq:E_log2}
  E_u(t)=(1+o(1))\frac{t}{(\log t)^2},\qquad
  B_u(t)=-(1+o(1))\frac{t}{(\log t)^2}.
\end{equation}
\end{enumerate}
In both regimes $|E_u(t)|$ and $|B_u(t)|$ exceed $t^{1/2+\eps}$ for every
fixed $\eps<1/2$ and all large $t$.
\end{proposition}

\begin{proof}
We use \cite{MontgomeryVaughan2007}: for some $c>0$,
\begin{equation}
\label{eq:pnt_inputs}
  \pi(y)=\li(y)+O\bigl(ye^{-c\sqrt{\log y}}\bigr),\quad
  M(y)\ll ye^{-c\sqrt{\log y}},\quad
  \mt(y):=\sum_{a\le y}\frac{\mu(a)}a\ll e^{-c\sqrt{\log y}},
\end{equation}
and $\sum_{n\ge1}\mu(n)\log n/n=-1$ (the value at $s=1$ of
$\zeta'(s)/\zeta(s)^2$). Hence
\begin{equation}
\label{eq:mtilde_integral}
  \int_1^\infty\mt(w)\frac{dw}w
  =\lim_{Y\to\infty}\Bigl(\mt(Y)\log Y-\sum_{a\le Y}\frac{\mu(a)\log a}a\Bigr)=1,
\end{equation}
with absolute convergence.

(i) Proposition~\ref{prop:fixed_R_corrected}, and $M(t)=o(t/\log t)$. For
$1<R<2$, $M=1$ on $[1,2)$.

(ii) Since $t<u^4$, $R<\sqrt t$, so $\log(t/a)\ge\frac12\log t$ for
$a\le R$. By \eqref{eq:prime_count_form},
\begin{equation}
\label{eq:E_three_terms}
  E_u(t)=-\sum_{a\le R}\mu(a)\li(t/a)
  -\sum_{a\le R}\mu(a)\bigl(\pi(t/a)-\li(t/a)\bigr)+M(R)\pi(u^2).
\end{equation}
The second term is $\ll t(1+\log R)e^{-c\sqrt{\log t/2}}=o(t/\log^2t)$. The
third is $\ll Re^{-c\sqrt{\log R}}\cdot u^2/\log u\ll e^{-c\sqrt{\log R}}\,
t/\log t$, which is $o(t/\log^2t)$ once $c\sqrt C>1$. For the first, with
$F(w)=w\li(t/w)$,
\begin{equation}
\label{eq:partial_summation_li}
  \sum_{a\le R}\mu(a)\li(t/a)=\mt(R)F(R)-\int_1^R\mt(w)F'(w)\,dw,
  \qquad
  F'(w)=\li(t/w)-\frac{t/w}{\log(t/w)}=\frac{t/w}{\log^2(t/w)}\Bigl(1+O\Bigl(\frac1{\log t}\Bigr)\Bigr).
\end{equation}
The boundary term is $\ll e^{-c\sqrt{\log R}}t/\log t=o(t/\log^2t)$. The
integral equals
$\frac t{\log^2t}\bigl(\int_1^R\mt(w)g_t(w)\frac{dw}w+O(1/\log t)\bigr)$ with
$g_t(w)=\log^2t/\log^2(t/w)\in[1,4]$ and $g_t\to1$ pointwise; by dominated
convergence and \eqref{eq:mtilde_integral} it is $(1+o(1))t/\log^2t$. Hence
$E_u(t)=(1+o(1))t/\log^2t$, and $B_u=M-E_u$ with $M(t)=o(t/\log^2t)$.
\end{proof}

\begin{remark}[Why $t/\log^2t$]
\label{rem:why_log2}
$\Cc(R)$ captures only the $1/\log t$ term, which tends to zero. The
$1/\log^2t$ term comes from $\zeta'(s)/\zeta(s)^2\to-1$, that is, from the
pole of $\zeta$ at $s=1$, and it does not decay in $R$. The version 5 claim
that the ``coherent fraction tends to zero'' is true for the first term and
does not describe the size of $E_u$.
\end{remark}

\begin{definition}[Remaining margin, version 5]
\label{def:remaining-margin}
For a target $T(t)=t^{1/2+\eps}$ put
\begin{equation}
\label{eq:margin-def}
  G_{u,\eps}(t)=\max\bigl(0,\ T(t)-|B_u(t)|\bigr),
\end{equation}
and call the large-prime part \emph{causally responsible} for a crossing of
$T$ if $|E_u(t)|\ge G_{u,\eps}(t)$ at the crossing.
\end{definition}

\begin{warning}[The margin is vacuous]
\label{warn:margin_vacuous}
With $B_u$ as in \eqref{eq:baseline_channel_def},
Proposition~\ref{prop:both_parts_large} gives $|B_u(t)|>T(t)$, hence
$G_{u,\eps}(t)=0$, for all large $t$ in both regimes. The test is always
passed and does not discriminate. ``Causal responsibility'' of one of two
cancelling terms of size $t/\log^kt$ for the value of their sum is not well
defined.
\end{warning}

\begin{warning}[Calibration remarks of version 5]
\label{warn:v5_calibration}
Version 5 required $|A(R)|\ll t^{-1/2+\eps}\log t$ for the large-prime part
to stay at the RH scale, and for $t=u^\alpha$ derived
$|A(R)|\ll R^{-\eta(\alpha,\eps)}$ with
$\eta=\alpha(1/2-\eps)/(\alpha-2)$ ($\eta=3/2,1,3/4,5/8$ at
$\alpha=3,4,6,10$). These remarks are \emph{superseded}: the coefficient is
$\Cc$, not $\Theta$ (Warning~\ref{warn:v5_main_term}); for $\alpha>4$ the
point $t$ lies outside the window of
Proposition~\ref{prop:linear-new-primes}, where
\eqref{eq:single-prime-copy} fails for $u^2<p\le\sqrt t$; and by
Proposition~\ref{prop:both_parts_large}, $E_u$ alone is never at the RH
scale, so only $B_u+E_u=M$ can be compared with $T$.
\end{warning}

\begin{remark}[First copy transitions]
\label{rem:first-copy-transitions}
For $1<R<2$ every copy has length $1$ and
$E_u(t)=-(\pi(t)-\pi(u^2))$, consistent with $\Cc(R)=1-1/R$. For $2<R<3$
copies of length $2$ appear but $M(2)=0$ ($\Cc=1/2$). For $3<R<4$ copies of
length $3$ enter with $M(3)=-1$, of opposite net sign.
\end{remark}

\section{Local large-prime decomposition of an excursion}
\label{sec:local_large_prime}

\noindent\emph{Status.} Proposition~\ref{prop:local_large_prime} and
Lemma~\ref{lem:inverse_log_sum} are original, with proofs; the input on
primes in short intervals is from the literature \cite{Huxley1972};
Heuristic~\ref{heu:local_drift} is heuristic.

\begin{proposition}[Local decomposition]
\label{prop:local_large_prime}
\emph{[unconditional]} Let $I_\nu=(u,v]$ be a complete excursion and
$y=\sqrt v$. For $u<t\le v$,
\begin{equation}
\label{eq:local_split}
  M(t)=\sum_{\substack{u<n\le t\\ \Pp(n)\le y}}\mu(n)
  -\sum_{a<y}\mu(a)\,\Pi_a(t),\qquad
  \Pi_a(t):=\#\Bigl\{p:\ \max\bigl(y,\tfrac ua\bigr)<p\le\tfrac ta\Bigr\},
\end{equation}
and in particular the bridge condition $M(v)=0$ reads
\begin{equation}
\label{eq:local_bridge}
  \sum_{\substack{u<n\le v\\ \Pp(n)\le\sqrt v}}\mu(n)
  =\sum_{a<\sqrt v}\mu(a)\,\Pi_a(v).
\end{equation}
\end{proposition}

\begin{proof}
$M(u)=0$. An $n\le v=y^2$ has at most one prime factor $p>y$; then
$n=ap$ with $a<v/y=y<p$, so $\mu(n)=-\mu(a)$, and $(a,p)\mapsto ap$ is a
bijection onto these $n$.
\end{proof}

\begin{remark}[What the large primes do inside an excursion]
\label{rem:local_meaning}
The decomposition holds for every excursion, with threshold $\sqrt{v_\nu}$
and no horizon $u_\nu^2$. The large-prime input consists of prime counts in
the intervals $(u/a,t/a]$, of length at most $\ell_\nu/a$ at height about
$v_\nu/a$. The asymptotic $\pi(x)-\pi(x-h)\sim h/\log x$ is known for
$x^\theta\le h\le x$ with fixed $\theta>7/12$ \cite{Huxley1972}, and not for
shorter intervals. Already for $a=1$, excursions with
$\ell_\nu<v_\nu^{7/12}$ have large-prime input without an unconditional
asymptotic.
\end{remark}

\begin{lemma}[Möbius-weighted inverse logarithms]
\label{lem:inverse_log_sum}
\emph{[unconditional]} With $C$ as in
Proposition~\ref{prop:both_parts_large}, uniformly for
$\exp(C(\log\log x)^2)\le z\le\sqrt x$,
\begin{equation}
\label{eq:inverse_log_sum}
  \sum_{a\le z}\frac{\mu(a)}{a\log(x/a)}=-\frac{1+o(1)}{(\log x)^2}.
\end{equation}
\end{lemma}

\begin{proof}
Partial summation gives
$\mt(z)/\log(x/z)-\int_1^z\mt(w)\,dw/(w\log^2(x/w))$; argue as in the proof
of Proposition~\ref{prop:both_parts_large}.
\end{proof}

\begin{heuristic}[Opposite drifts]
\label{heu:local_drift}
This is a heuristic argument. Let $u\ge\sqrt v\exp(C(\log\log v)^2)$. Near
$t\in I_\nu$ the counts $\Pi_a$ grow for $a<t/\sqrt v$, at rate
$1/(a\log(t/a))$ if primes near $t/a$ have their expected local density
uniformly in $a$ (beyond what is known). Ignoring these difficulties, over
$(t,t+h]$ the large-prime part changes by about
$-h\sum_{a<t/\sqrt v}\mu(a)/(a\log(t/a))=(1+o(1))h/\log^2t$ by
Lemma~\ref{lem:inverse_log_sum}. Since $M(v)-M(u)=0$, the
$\sqrt v$-smooth part drifts at rate about $-1/\log^2v$. The drift does not
depend on $\sigma_\nu$. For $\ell_\nu\gg(\log v_\nu)^4$ both parts exceed
$\sqrt{\ell_\nu}$, and \eqref{eq:local_bridge} is a balance between two
opposite drifts: the local analogue of
Proposition~\ref{prop:both_parts_large}.
\end{heuristic}

\section{A model for the black curve}
\label{sec:model}

\noindent\emph{Status.} Heuristic throughout.

\begin{heuristic}[Scale and period]
\label{heu:model}
This is a heuristic argument. Motivated by \eqref{eq:explicitM}, one is led
to model $M(x)\approx\sqrt x\,F(\log x)$ with $F$ almost periodic, dominated
at low frequency by $\gamma_1=14.1347\ldots$. Within an excursion of length
$\asymp x$ this suggests $H_\nu\asymp\sqrt x$ and
$\sqrt{H_\nu}\ell_\nu\asymp x^{5/4}$, so the black curve should oscillate at
scale
\begin{equation}
\label{eq:model_scale}
  |\Anet(x)|\asymp x^{5/4},
\end{equation}
with a period $2\pi/\gamma_1$ in $\log x$, that is, a factor
$e^{2\pi/\gamma_1}\approx1.56$ in $x$.
\end{heuristic}

\section{Numerical data}
\label{sec:data}

\noindent\emph{Status.} Observations recorded in version 5; no new
computation was carried out for version 6.

\begin{experiment}[Excursion census]
\label{exp:census}
\emph{Object:} the excursion data $(u_\nu,\tau_\nu,v_\nu,H_\nu,\sigma_\nu)$
and the curves $\Ap,\Am,\Anet$.
\emph{Range:} $n\le N=455\,052\,511$, with $Z(N)=87\,385$ excursions.
\emph{Expected outcome:} the bound \eqref{eq:transfer_bound}; the prediction
$Z(N)N^{-1/4}\approx6.0\times10^2$ excursions longer than
$\sqrt N\approx2.13\times10^4$ from Heuristic~\ref{heu:renewal}.
\emph{Observed outcome (version 5):}
\begin{equation}
\label{eq:census_values}
  \Ap(N)=6.3449\times10^9,\quad \Am(N)=-7.0877\times10^9,\quad
  \Anet(N)=-7.4279\times10^8,
\end{equation}
\begin{equation}
\label{eq:census_max}
  \max_\nu\frac{H_\nu}{\sqrt{\ell_\nu}}=2.0743,\qquad
  \max_\nu\frac{\beta_\nu}{\ell_\nu^{3/2}}=1.0214 .
\end{equation}
The count of excursions longer than $\sqrt N$ and the maxima restricted to
$\ell_\nu>\sqrt N$ were not recorded. At this range $N^{0.05}\approx2.7$, so
factors $x^\eps$ are not separable from constants.
\end{experiment}

\begin{warning}[Constant in the transfer bound]
\label{warn:census_c}
Version 5 compared $|\Anet(N)|$ with \eqref{eq:transfer_bound} using
$c=\max_{10^3\le n\le N}|M(n)|n^{-1/2}$, reporting a bound about $30$ times
larger than $|\Anet(N)|$. That $c$ does not satisfy the hypothesis of
Lemma~\ref{lem:transfer}, which requires all $n\le x^*$; since $M(1)=1$, any
admissible $c$ is $\ge1$. The comparison must either use the full maximum or
split off the range below $10^3$.
\end{warning}

\section{Open questions and future work}
\label{sec:open_questions}

\noindent\emph{Status.} Open problems and proposed experiments; none has
been carried out.

\begin{enumerate}[label=\textup{(\arabic*)}]
\item Question~\ref{q:finiteness_expanding}: finiteness of
information-expanding excursions.
\item Whether \eqref{eq:rh_area} is equivalent to RH
(Remark~\ref{rem:onedirection}).
\item Whether the pruned local condition \eqref{eq:pruned_Q} holds, and a
third bridge constraint (Remark~\ref{rem:conditioning}).
\item Experiment: the number of excursions with $\ell_\nu>\sqrt N$ against the
prediction $\approx600$, and the maxima of $H_\nu/\sqrt{\ell_\nu}$ and
$\beta_\nu/\ell_\nu^{3/2}$ over those excursions.
\item Experiment: split each excursion by \eqref{eq:local_split} and test
Heuristic~\ref{heu:local_drift}, in particular the independence of the drift
from $\sigma_\nu$.
\item Experiment: frequency of touching runs (Definition~\ref{def:crossing_touching})
and of ballistic excursions (Heuristic~\ref{heu:local_fails}).
\end{enumerate}

\section*{Notebook History}
\label{sec:history}

\noindent\textbf{Version 5.} Single section: definitions, exact formulas,
internal correlations, bounds, criteria, prime horizon, copy propagation,
model, census data.

\medskip
\noindent\textbf{Version 6.} Restructured into labelled sections following
the authoring rules. Corrections: null excursions and touching zeros
(Definitions~\ref{def:null}, \ref{def:crossing_touching},
Lemma~\ref{lem:critical}); Remark~\ref{rem:omega};
Remark~\ref{rem:onedirection}; hypothesis of Lemma~\ref{lem:transfer};
constant in Remark~\ref{rem:chowla-scale}; $j=1$ in
Proposition~\ref{prop:double-exp}; \eqref{eq:Ny-fixed-R}; main term
(Proposition~\ref{prop:fixed_R_corrected}, Warning~\ref{warn:v5_main_term}).
New: Warning~\ref{warn:information_language}, Lemma~\ref{lem:rough_part},
Proposition~\ref{prop:prime_count_form},
Proposition~\ref{prop:both_parts_large},
Section~\ref{sec:local_large_prime}, Heuristic~\ref{heu:local_fails},
Corollary~\ref{cor:pruned_local}, proofs of
Propositions~\ref{prop:L1} and~\ref{prop:dipoles}. Archived as superseded:
the calibration remarks (Warning~\ref{warn:v5_calibration}) and the margin
test (Warning~\ref{warn:margin_vacuous}).

\begin{thebibliography}{99}

\bibitem{Huxley1972}
M.~N. Huxley, On the difference between consecutive primes,
\emph{Invent. Math.} \textbf{15} (1972), 164--170.
doi:10.1007/BF01418933.

\bibitem{KotnikTeRiele2006}
T.~Kotnik and H.~te Riele, The Mertens conjecture revisited, in
\emph{Algorithmic Number Theory (ANTS VII)}, Lecture Notes in Comput. Sci.
\textbf{4076}, Springer, Berlin, 2006. doi:10.1007/11792086\_12.

\bibitem{MontgomeryVaughan2007}
H.~L. Montgomery and R.~C. Vaughan, \emph{Multiplicative Number Theory I:
Classical Theory}, Cambridge Stud. Adv. Math. \textbf{97}, Cambridge
University Press, 2007. ISBN 978-0-521-84903-6.

\bibitem{Ng2004}
N.~Ng, The distribution of the summatory function of the Möbius function,
\emph{Proc. London Math. Soc. (3)} \textbf{89} (2004), no.~2, 361--389.
doi:10.1112/S0024611504014741. arXiv:math/0310381.

\bibitem{OdlyzkoTeRiele1985}
A.~M. Odlyzko and H.~J.~J. te Riele, Disproof of the Mertens conjecture,
\emph{J. Reine Angew. Math.} \textbf{357} (1985), 138--160. MR0783538.

\bibitem{Titchmarsh1986}
E.~C. Titchmarsh, \emph{The Theory of the Riemann Zeta-Function}, 2nd ed.,
revised by D.~R. Heath-Brown, Clarendon Press, Oxford, 1986.

\bibitem{Walfisz1963}
A.~Walfisz, \emph{Weylsche Exponentialsummen in der neueren Zahlentheorie},
VEB Deutscher Verlag der Wissenschaften, Berlin, 1963.

\end{thebibliography}

\end{document}
